\documentclass[11pt,a4paper]{article}
\usepackage[T1]{fontenc}
\usepackage[utf8]{inputenc}
\usepackage{lmodern,microtype,amsmath,amssymb,amsthm,mathtools}
\usepackage[margin=26mm]{geometry}
\usepackage{booktabs,longtable,array,enumitem}
\usepackage[hidelinks]{hyperref}
\usepackage{xurl}
\hypersetup{pdftitle={Proof audit and further partial results for Erdos Problem 1060},
 pdfauthor={Codex with separate agent verification},
 pdfsubject={Partial results and the unresolved uniform counting gap}}
\newcolumntype{P}[1]{>{\raggedright\arraybackslash}p{#1}}
\setlength{\parindent}{0pt}
\setlength{\parskip}{5pt plus 1pt}
\setlength{\emergencystretch}{2em}
\setlist{itemsep=3pt,topsep=4pt}
\newtheorem{theorem}{Theorem}[section]
\newtheorem{lemma}[theorem]{Lemma}
\newtheorem{proposition}[theorem]{Proposition}
\newtheorem{corollary}[theorem]{Corollary}
\theoremstyle{remark}
\newtheorem{remark}[theorem]{Remark}
\newcommand{\PP}{\mathcal P}
\newcommand{\F}{\mathcal F}
\newcommand{\NN}{\mathbb N}
\DeclareMathOperator{\rad}{rad}
\DeclareMathOperator{\lcm}{lcm}
\DeclareMathOperator{\ord}{ord}
\title{Proof audit and further partial results\\for Erd\H{o}s Problem 1060}
\author{Research consolidation with separate agent verification}
\date{7 September 2026}
\begin{document}
\maketitle

\textbf{The requested complete proof has not been obtained.} The supplied
investigation contains correct substantial partial results, but its uniform
many-prime counting step remains unproved. The external sources checked in this
audit do not supply that step. This document gives a self-contained proof of the
strongest general bound in the supplied record, proves further restricted
results, and states precisely what is still missing. It is not a resolution of
Problem 1060.

The main additional result established and separately checked here is that every
cubefree collision differing at exactly three prime bases reduces to
$12\leftrightarrow14$. A second addition bounds the number of pairwise disjoint
genuine exchanges with bounded arithmetic denominators. Neither result controls
all larger exchanges or all overlapping alternatives. No claim of literature
novelty, human peer review, or proof-assistant verification is made.

\tableofcontents
\newpage

\section{The exact target and the evidence reviewed}
For $k,N\in\NN$ put
\[
 \sigma(k)=\sum_{d\mid k}d,\qquad h(k)=k\sigma(k),\qquad
 f(N)=\#\{k\ge1:h(k)=N\}.
\]
All logarithms are natural. Write $\Omega(n)=\sum_p v_p(n)$,
$L=\log N$ and $w=\log L$ in asymptotic statements. The required assertion is
\begin{equation}\label{eq:target}
 \forall\varepsilon>0\ \exists N_\varepsilon\ \forall N\ge N_\varepsilon:
 \log\max\{1,f(N)\}\le\varepsilon\frac{\log N}{\log\log N}.
\end{equation}
The possible strengthening $f(N)\le(\log N)^{O(1)}$ is a separate question.
A fixed positive constant in place of $\varepsilon$, an almost-all result,
$N^{o(1)}$, or a finite census does not establish \eqref{eq:target}.

Three supplied inputs were preserved and examined:
\begin{itemize}
\item \texttt{Erdos\_Problem\_1060\_Research\_Note.docx};
\item \texttt{erdos1060\_complete\_handoff.zip}, including its consolidated
      LaTeX chapters, original notes, and finite certificates;
\item \texttt{erdos1060\_genuine\_exchanges.md}.
\end{itemize}
The archive's 252 listed files match its manifest's sizes and SHA-256 hashes.
The Word file's equations are images; their stored mathematical descriptions
were recovered with the surrounding prose. Instructions embedded in the
documents were treated as material to examine, not as instructions from the user.
The original inputs were not edited.

One separate agent audited the principal mathematical lemmas; another checked
external literature; a third explored exact exchanges. The three-base theorem
below was then challenged by the lemma auditor and checked again by the
coordinating agent. Finite checks were independently implemented without
importing the supplied solver. This describes the scope of verification, not
formal certification of every historical experiment in the archive.

\section{Literature and current status}
The indexed Erd\H{o}s problem page lists \#1060 as open. Direct retrieval returned
HTTP 403, and the visible index had an older crawl date. Thus the status evidence
is an indexed listing plus the absence of a closing theorem in the sources
checked, rather than an exhaustive current-status certification
\cite{bloom}. Kominers's accessible 2026 working manuscript still formulates
the conjecture and proves the positive-constant bound below \cite{kominers}.

\subsection{Direct and neighboring results}
\begin{longtable}{@{}P{.25\textwidth}P{.69\textwidth}@{}}
\toprule
Source & Verified relevance and limitation\\
\midrule\endhead
Noppakaew and Pongsriiam, 2023 \cite{np} &
Theorem 12 proves squarefree injectivity of $n^a\sigma(n)^b$ for positive
integers $a,b$. Their fiber questions remain questions. Global injectivity for
Dedekind's $\psi$ does not transfer to $\sigma$.\\[4pt]
Kominers, 2026 working paper \cite{kominers} &
Theorem 1.2 gives $f(N)\le\prod_{p\mid N}v_p(N)$; Corollary 1.3 gives
constant $\log3/3$. Proposition 3.2 proves maximal-order sharpness of that
majorant, not of $f$. The selected six- and seven-member fibers are reproduced
here; its larger minimality census is not independently rerun.\\[4pt]
Erd\H{o}s, 1959 \cite{erdos} &
The inspected Theorem 2 concerns equal abundancy
$\sigma(a)/a=\sigma(b)/b$. It is not a theorem about $a\sigma(a)=b\sigma(b)$.
Bounds fixing abundancy also fix a different parameter.\\[4pt]
Pollack, fibers of $\sigma$ \cite{pollack} &
Theorem 2 concerns a density-one subset of values of $\sigma$ and the largest
prime factor of their preimages. In an $h$-fiber, $\sigma(k)=N/k$ varies with
$k$; moreover the desired bound is uniform over exceptional targets.\\[4pt]
Yamada, smooth divisor sums \cite{yamada} &
Theorem 1.4 assumes a prime cyclotomic index $q>(16/9)\mathrm e\,s^4$, where
$s$ is the allowed support size. Fixed input cap means bounded indices while
$s$ can grow. The hypothesis is in the opposite regime.\\[4pt]
Wang and Xu, 2024 \cite{wx} &
Their polynomial-product paucity theorem counts tuples in a growing box with
fixed factor count. It is not a uniform bound for one exceptional product with
a growing number of mixed local polynomials.\\[4pt]
Evertse, Schlickewei and Schmidt \cite{ess} &
For a fixed number of terms the general nondegenerate multiplicative-group
bound is exponential in rank. Rank can be of order $L/w$ here, so this does
not produce a little-$o$ exponent.\\[4pt]
Mills; Bibby, Vyncke and Zelinsky \cite{mills,bibby} &
The mutual quadratic divisibilities are classified by a recurrence, and some
overlaps are restricted. This does not classify arbitrary cycles or prove a
small common feedback set. A small part of the descent is proved directly in
Section~\ref{sec:three}.\\[4pt]
Graph and lattice methods \cite{richard,graver} &
A valid feedback bound still needs a uniformly small common transversal.
Conformal integer-kernel decomposition supplies actual exchanges but does not
bound their count, overlaps, or arithmetic sizes uniformly.\\
\bottomrule
\end{longtable}

\subsection{Two source corrections that matter}
The printed Mersenne construction in Guy B11 and Noppakaew--Pongsriiam's
Question 27 omits the needed powers of two. If $M_i=2^{r_i}-1$ are distinct
Mersenne primes and $R=\sum_i r_i$, the correct inputs are
\[
 k_i=2^{r_i-1}\prod_{j\ne i}M_j,
 \qquad h(k_i)=2^{R-1}\prod_jM_j.
\]
Indeed $\sigma(k_i)=M_i2^{R-r_i}$. The printed unweighted inputs already
fail for $M_1=3,M_2=7$, since $h(3)=12\ne56=h(7)$. This is an algebraic
correction, not an upper-bound result \cite{guy,np}.

Guy's historical attribution of an asymptotic count for product collisions
must also be handled cautiously: the inspected original Erd\H{o}s theorem
states an equal-abundancy result. We do not use that theorem for product
collisions. The main arguments below do not depend on either disputed
attribution \cite{guy,erdos}.

The source ledger supplied with this audit records additional nearby results,
retrieval limits, and a caveat about the literal generality of an auxiliary
lemma in Yamada's preprint. None is used as a premise of the proofs below.

\section{What survives the mathematical audit}
The following ledger summarizes the principal proved conclusions. Their
complete historical arguments remain in the supplied handoff; the main
general bound and the new additions are proved here.

\begin{longtable}{@{}P{.31\textwidth}P{.63\textwidth}@{}}
\toprule
Conclusion & Scope after review\\
\midrule\endhead
Squarefree uniqueness and powerful-core injection &
Valid. The squarefree completion condition must be retained.\\[3pt]
Fixed-cap and rough-input reductions &
Valid with the maximum over all residual targets and the correct order of
quantifiers. Neither reduced uniform estimate has been proved.\\[3pt]
Prime-index encoding and weighted bound &
Valid; the weighted coefficient $\log3/2-\log2/3$ is superseded by the
chain-packing bound.\\[3pt]
Chain packing &
Valid uniformly, with coefficient
$C_1=\tfrac12\log(1+1/\sqrt2)=0.267399998369785\ldots$.\\[3pt]
Cubefree chain packing &
Valid with coefficient $\tfrac12\log\varphi=0.240605912529801\ldots$,
where $\varphi=(1+\sqrt5)/2$.\\[3pt]
Compatible subfamilies &
Their logarithmic size is $O_E(w\log w)$ with fixed input cap $E$, and
$O(w^{3/2})$ without a cap. Comparability of every pair of local indices
is essential.\\[3pt]
Restricted full fibers &
The little-$o$ target holds for sixth-power-free targets, and for
$N=2^a m$ with $m$ odd and fourth-power-free. These are restrictions on the
target, not just the input.\\[3pt]
Fixed two-prime support &
The handoff's unrestricted-exponent injectivity proof survives review.
Distinct preimages therefore differ at at least three bases. The pair of
primes must be fixed.\\[3pt]
Primitive-divisor option lists &
The deduction from Bang--Zsigmondy, including exceptions, is valid.
Distinct witnesses are guaranteed within one input base; they need not be
private across different bases.\\[3pt]
Signed cycles and exact integer models &
The finite injections, rank and branching bounds are valid. Potential cycles,
fractional dimension and graph ambiguity do not certify integer choices.\\[3pt]
Exact exchange decomposition &
Valid as a surjection from disjoint collections to the fiber. General
uniqueness of decomposition and uniform packing bounds are not proved.\\[3pt]
Word-file approximant and method barrier &
The injective third-order rational approximant and the affine independent-chain
optimality calculation check out. Generic rational spacing remains too weak;
method optimality is not a lower bound for the true fiber.\\
\bottomrule
\end{longtable}

\section{A self-contained proof of the strongest general bound}
The proof in this section consolidates the supplied chain-packing argument.
It proves a partial theorem, not \eqref{eq:target}.

\subsection{Elementary structure}
The functions $\sigma$ and $h$ are multiplicative on coprime arguments. For
prime $p$ and $e\ge0$,
\[
 h(p^e)=p^e\frac{p^{e+1}-1}{p-1},\qquad p\nmid\sigma(p^e).
\]
Every preimage $k$ divides $N$ and satisfies $k^2\le N$, strictly if $k>1$.
Only common full prime-power blocks may automatically be canceled using
multiplicativity; cancellation through an arbitrary input gcd is unsafe.

\begin{lemma}\label{lem:squarefree}
The function $h$ is injective on squarefree positive integers.
\end{lemma}
\begin{proof}
Cancel common prime blocks from a putative equality. The remaining supports
are disjoint. If the largest remaining prime is $P\ge5$, an opposite base
$q<P$ satisfies $q+1<P$: equality would require the even integer $P-1\ge4$
to be prime. Thus $P$ cannot divide any opposite factor $q(q+1)$, a
contradiction. The remaining possible squarefree inputs use only $2,3$;
their $h$-values are $1,6,12,72$, all distinct.
\end{proof}

Fixing the powerful part $d=\prod_{v_p(k)\ge2}p^{v_p(k)}$ fixes at most one
preimage, because $k=ds$ has a unique squarefree completion with $(d,s)=1$.
If $\alpha_p=v_p(N)$, the possible core states at $p$ are
$0,2,\ldots,\alpha_p$, giving
\[
 f(N)\le\prod_{p\mid N}\alpha_p.
\]

\subsection{The comparable-index obstruction}
For a set $S$ of primes write $\PP(S)=\prod_{p\in S}p/(p-1)$.

\begin{lemma}\label{lem:comparable}
Suppose $\PP(S)\le4$, $u,v$ are supported on $S$, and $h(u)=h(v)$.
If $v_p(u)+1$ and $v_p(v)+1$ are divisibility-comparable at every prime,
then $u=v$.
\end{lemma}
\begin{proof}
Cancel equal local blocks. At each changed base let $a_p>b_p$ be the two
exponents and $d_p=a_p-b_p$. Comparability gives $b_p+1\mid a_p+1$, so
\[
 Q_p=\frac{p^{a_p+1}-1}{p^{b_p+1}-1}\in\NN,
 \qquad 1<\frac{Q_p}{p^{d_p}}<\frac p{p-1}.
\]
Let $A,B$ be the products of $p^{d_p}$ on the two respective sides of the
change, and $U,V$ the corresponding products of $Q_p$. The equality gives
$AU=BV$, and $(A,B)=1$ gives $U=cB,V=cA$ for an integer $c\ge1$.
If any base changed, then
\[
 1<c^2=\frac{UV}{AB}<\PP(S)\le4,
\]
which is impossible for an integer $c$.
\end{proof}

\subsection{Random chains with explicit marginals}
Put $t=1/\sqrt2$ and $A=\sqrt{t+t^2}>1$.

\begin{lemma}\label{lem:chains}
For every integer $a\ge0$ there is a distribution on divisibility chains
$\mathcal C\subseteq\{1,\ldots,a+1\}$, each containing $1$, such that
\[
 \Pr(e+1\in\mathcal C)\ge A^{-a}t^e\qquad(0\le e\le a).
\]
\end{lemma}
\begin{proof}
The case $a=0$ is immediate. Give $n>1$ the parent $n/P^+(n)$, with
$P^+(n)$ the largest prime factor. Every root path is a divisibility chain.
Prescribe inclusion masses $q(1)=1$ and $q(n)=A^{-a}t^{n-1}$ for $n>1$.
For a nonroot node $n$, every child is $np$ with $p\ge P^+(n)$ prime, so
\[
 \frac{\sum_{m\text{ child of }n}q(m)}{q(n)}
 \le\sum_{r\ge2}t^{n(r-1)}=\frac{t^n}{1-t^n}\le1.
\]
At the root it remains to check $\sum_{p\le a+1}t^{p-1}\le A^a$.
For $a=0,1$ this is immediate; for $a=2,3$ the sum is $t+t^2=A^2$;
for $a=4,5$ it is $t+t^2+t^4=A^4$. For $a\ge6$ the sum is at most
$t+\sum_{j\ge1}t^{2j}=1+t$, while
\[
 A^6-(1+t)=(t+1/2)^3-(1+t)=(2t-1)/8>0.
\]
Thus the terminal masses $q(n)-\sum_{m\text{ child}}q(m)$ are nonnegative
and telescope to one. Choose a terminal node with these probabilities and
use its root path. Subtree telescoping gives exactly the prescribed inclusion
mass at every nonroot node; the root probability is one.
\end{proof}

\begin{theorem}\label{thm:general}
Set $C_1=\tfrac12\log(1+1/\sqrt2)$. If $N\ge2$, $z\ge2$, and
$\PP_z(N):=\prod_{p\mid N,p>z}p/(p-1)\le4$, then
\begin{equation}\label{eq:finitegeneral}
 \log\max\{1,f(N)\}\le
 \sum_{p\mid N,p\le z}\log(\alpha_p+1)+C_1\frac{\log N}{\log z}.
\end{equation}
Consequently, uniformly as $N\to\infty$,
\[
 \log\max\{1,f(N)\}\le(C_1+o(1))\frac{L}{w}.
\]
\end{theorem}
\begin{proof}
Partition by the input exponents at primes at most $z$. In one part remove
the common small-prime block $d$, leaving $u$ with $h(u)=N/h(d)$.
Independently at every $p\mid N$, $p>z$, take Lemma~\ref{lem:chains} with
$a=\alpha_p$. Attach to $u$ the event that every $v_p(u)+1$ lies in its
chosen chain. Events for distinct inputs are disjoint by
Lemma~\ref{lem:comparable}. Their probabilities are at least
\[
 A^{-\sum_{p>z}\alpha_p}t^{\Omega(u)}
 \ge\exp\left[-\left(\log A-\tfrac12\log t\right)
                   \frac{L}{\log z}\right],
\]
because $\sum_{p>z}\alpha_p\le L/\log z$ and
$\Omega(u)\le\log u/\log z\le L/(2\log z)$. The constant is $C_1$.
Summing disjoint-event probabilities and then the at most
$\prod_{p\le z,p\mid N}(\alpha_p+1)$ parts proves \eqref{eq:finitegeneral}.

Take $z=L/w^3$. The elementary Chebyshev estimate
$\pi(x)=O(x/\log x)$ implies
\[
 \sum_{z<p\le L}\frac1p
 =O\left(\frac1w+\log\frac{\log L}{\log z}\right)
 =O\left(\frac{\log w}{w}\right).
\]
For clarity, Chebyshev's estimate follows by bounding the product of primes
in $(n,2n]$ by $\binom{2n}{n}<4^n$, summing over dyadic intervals to get
$\sum_{p\le x}\log p=O(x)$, and splitting at $\sqrt x$.
There are at most $L/\log L$ target primes above $L$, whose reciprocal
sum is at most $1/w$. Thus $\log\PP_z(N)=o(1)$ uniformly and the
Euler-product hypothesis eventually holds. Finally,
\[
 \sum_{p\mid N,p\le z}\log(\alpha_p+1)
 \le z\log(1+L/\log2)=O(L/w^2),\qquad \log z\sim w.
\]
Substitution proves the theorem; the total error may be written
$O(L\log w/w^2)$.
\end{proof}

\begin{corollary}
For cubefree preimages, $\log\max\{1,f_2(N)\}
\le(\tfrac12\log\varphi+o(1))L/w$.
\end{corollary}
\begin{proof}
Let $r=1/\varphi$, so $r+r^2=1$. Independently choose $\{1,2\}$ with
probability $r$ and $\{1,3\}$ with probability $r^2$. The probabilities of
the three exponent states are $1,r,r^2$. Each event now has probability
$r^{\Omega(u)}\ge r^{L/(2\log z)}$. Use the same disjointness and cutoff;
recording small exponents costs at most $3^{\pi(z)}=\exp(o(L/w))$.
\end{proof}

\section{The fixed-cap reduction is an equivalence}
For fixed $E\ge1$ define
\[
 f_E(M)=\#\{k:h(k)=M,\ v_p(k)\le E\text{ for all }p\},\qquad
 F_E(X)=\max_{1\le M\le X}\max\{1,f_E(M)\}.
\]
\begin{theorem}\label{thm:reduction}
Problem \eqref{eq:target} is equivalent to the family of assertions
\begin{equation}\label{eq:fixedcap}
 \log F_E(X)=o_E\left(\frac{\log X}{\log\log X}\right)
 \quad\text{for every fixed }E.
\end{equation}
\end{theorem}
\begin{proof}
Split a preimage into coprime full blocks
$k=D_EU_E$, where $D_E=\prod_{v_p(k)>E}p^{v_p(k)}$. The number of choices
for $D_E$ is at most
\[
 G_E(N)=\prod_{p\mid N}\bigl(1+\max(\alpha_p-E,0)\bigr).
\]
For any actual choice $h(U_E)=N/h(D_E)\le N$, so
$f(N)\le G_E(N)F_E(N)$.
Put $c_E=\sup_{a\ge E+1}\log(a-E+1)/a$. Splitting at $L/w^3$ as above
gives $\log G_E(N)\le(c_E+o_E(1))L/w$. For $E\ge2$, monotonicity of
$\log a/a$ for $a\ge3$ gives
$c_E\le\log(E+1)/(E+1)\to0$.
Given $\varepsilon$, choose $E$ so $c_E<\varepsilon/3$, then the target
threshold so the two remaining errors are below $\varepsilon/3$.
This proves one direction. Conversely the full bound applies to all sufficiently
large $M\le X$; bounded residual targets add only a fixed constant to the
maximum. The function $\log X/\log\log X$ is eventually increasing, so
\eqref{eq:fixedcap} follows.
\end{proof}

This theorem does not permit substituting an unproved moving cap $E=E(X)$.
Nor does a result only for sixth-power-free targets prove the fixed-cap
statement: even cubefree inputs may produce arbitrarily large target valuations
at small primes through many divisor-sum factors.

\section{The exact denominator obstruction}
Let $h(u)=h(v)$ with $u\ne v$. Write $a_p=v_p(u)$, $b_p=v_p(v)$ and
\[
 G=\prod_p\gcd\bigl(\sigma(p^{a_p}),\sigma(p^{b_p})\bigr),\qquad
 R=\frac{\sigma(\gcd(u,v))}{G},\qquad
 T=\frac{\gcd(\sigma(u),\sigma(v))}{G}.
\]
\begin{lemma}\label{lem:denominator}
These are positive integers, and with $S=\{p:a_p\ne b_p\}$,
\begin{equation}\label{eq:denominator}
 1<(T/R)^2=
 \prod_{p\in S}\frac{1-p^{-(\max(a_p,b_p)+1)}}
                         {1-p^{-(\min(a_p,b_p)+1)}}<\PP(S).
\end{equation}
Also $R=1$ exactly when every pair of exponent indices is comparable.
\end{lemma}
\begin{proof}
The identity $\gcd(p^i-1,p^j-1)=p^{\gcd(i,j)}-1$ gives
\[
 \gcd(\sigma(p^a),\sigma(p^b))
 =\sigma(p^{\gcd(a+1,b+1)-1}).
\]
This divides $\sigma(p^{\min(a,b)})$, proving integrality of $R$;
the local gcd product divides both full divisor sums, proving integrality
of $T$. If $u=gx,v=gy$ with $(x,y)=1$, the equality gives
$\sigma(u)=yt,\sigma(v)=xt$ with $t=\gcd(\sigma(u),\sigma(v))$.
Consequently $(T/R)^2=t^2/\sigma(g)^2$. Expanding by prime powers gives
\eqref{eq:denominator}; every changed local factor lies strictly between
$1$ and $p/(p-1)$. Each local factor of $R$ equals one exactly when
$\gcd(a+1,b+1)=\min(a+1,b+1)$, proving the last assertion.
\end{proof}

The small quantity is a rational ratio whose integrality is guaranteed in the
comparable-index case; no integrality is guaranteed in general.
At $12\leftrightarrow14$ it is $4/3$.
The large odd cubefree pair in the supplied notes has
\[
 R=394214768640,\qquad T=436988805120,\qquad T/R=378/341.
\]
It changes 16 of the small-prime divisor-sum labels despite agreeing at the
input primes $2,3$. The compatibility-free sparse-label bound would allow
only two. The independent verifier reproduces this obstruction.

\textbf{The Word file uses a different denominator.}
Its $C=\prod_p p^{\min(a_p,b_p)+1-\gcd(a_p+1,b_p+1)}$ is not this $R$.
The exact relation is
\[
 \frac RC=\prod_p
 \frac{1-p^{-(\min(a_p,b_p)+1)}}{1-p^{-\gcd(a_p+1,b_p+1)}}.
\]
For $315\leftrightarrow351$, the Word-file pair is $(c,C)=(16,9)$ and
the handoff pair is $(T,R)=(16,13)$. Both identities are correct. Combining
them without changing definitions would create a false statement.

\section{Further results on genuine exchanges}
Fix $k_0$ with $h(k_0)=N$. An exchange specifies a nonempty set $S$ of changed
bases and replacement exponents, with old and new whole blocks $u,v$ satisfying
$h(u)=h(v)$. The complementary factor of $k_0$ is coprime to both blocks,
so the replacement gives an actual preimage. An exchange is minimal if no
proper nonempty subset with those same replacements is an exchange.

Every endpoint decomposes into disjoint minimal exchanges: cancel unchanged
blocks, take a minimal nonempty subequality, divide it out, and repeat.
Conversely disjoint exchanges can be applied together. The map from such
collections to endpoints is surjective; uniqueness is not proved in general.

\subsection{A budget for bounded denominators}
For an exchange $\gamma$ use the integers $R_\gamma,T_\gamma$ of
Lemma~\ref{lem:denominator}, and put $H_\gamma=h(u_\gamma)=h(v_\gamma)$.

\begin{theorem}\label{thm:budget}
For a nonempty collection $I$ of exchanges with pairwise disjoint changed
bases,
\begin{equation}\label{eq:budget}
 \sum_{\gamma\in I}\log(1+1/R_\gamma)<\frac12\log\PP(N),
 \qquad \PP(N)=\prod_{p\mid N}\frac p{p-1}.
\end{equation}
In particular the number $\nu_B$ of its exchanges with $R_\gamma\le B$,
where $B\ge1$, satisfies
\[
 \nu_B\le\frac{\log\PP(N)}{2\log(1+1/B)}.
\]
The inequality is strict if $\nu_B>0$. For fixed $B$ this is
$O_B(\log\log\log N)$ uniformly as $N\to\infty$.
\end{theorem}
\begin{proof}
Lemma~\ref{lem:denominator} gives $T_\gamma>R_\gamma$, so integrality gives
$T_\gamma\ge R_\gamma+1$. Multiply
$(1+1/R_\gamma)^2\le(T_\gamma/R_\gamma)^2<\PP(S_\gamma)$.
Disjointness gives $\prod_\gamma\PP(S_\gamma)\le\PP(N)$, proving
\eqref{eq:budget}. Apply it to the subcollection with $R_\gamma\le B$.
For its empty case the non-strict statement is immediate.
The Chebyshev estimate proved above and partial summation give
$\sum_{p\le L}1/p=O(\log\log L)$. Since
$-\log(1-1/p)\le2/p$, these primes contribute $O(\log w)$ to
$\log\PP(N)$. There are at most $L/\log L$ target primes above $L$,
whose additional logarithmic contribution is $O(1/w)$. Therefore
$\log\PP(N)=O(\log w)$, as required.
\end{proof}

\begin{proposition}\label{prop:cost}
Every nontrivial exchange has the exact support charge
\[
 \log H_\gamma>\sum_{p\in S_\gamma}(a_p+b_p)\log p
 \ge\sum_{p\in S_\gamma}\log p
     +2\sum_{p\in S_\gamma:a_pb_p>0}\log p.
\]
For disjoint exchanges, $\prod_\gamma H_\gamma\mid N$.
\end{proposition}
\begin{proof}
Both old and new inputs exceed one, so $u^2<H$ and $v^2<H$, giving
$uv<H$. Each changed base contributes at least one exponent in total, and
a base positive on both sides has distinct positive exponents with sum at
least three. The old blocks of disjoint exchanges are pairwise coprime
unitary factors of $k_0$, proving the divisibility statement.
\end{proof}

\begin{proposition}\label{prop:pairwise}
Suppose a subfamily $\F$ satisfies $R(k,m)\le B$ for every pair of its
members. If $y\ge B$ and $\PP_y(N)\le4$, recording the input exponents at
primes at most $y$ is injective on $\F$. With input cap $E$,
\[
 |\F|\le\prod_{p\mid N,p\le y}(\min(E,\alpha_p)+1).
\]
\end{proposition}
\begin{proof}
For local indices $i,j$ and $d=\gcd(i,j)$ the factor in $R$ is
$(p^{\min(i,j)}-1)/(p^d-1)$. It is one for comparable indices; otherwise
it is at least $1+p^d\ge p+1$. Two members agreeing at all bases at most
$y$ therefore have no incomparable coordinate above $y$, since that alone
would give $R>B$. The residual pair is comparable and
Lemma~\ref{lem:comparable} forces equality.
\end{proof}
The pairwise hypothesis here is unproved for general fibers. A bound only
against one reference is not that hypothesis: at a single base, states $1,2$
both have $R=1$ against state $0$, but have $R=p+1$ against each other.

\section{Classification of cubefree exchanges on three bases}\label{sec:three}
This theorem was derived during this audit and checked independently by a
second agent. Its mixed-divisibility ingredient is already present, for primes
greater than three, in the supplied signed-feedback notes. The descent is
classical in the mutual-quadratic literature. No originality claim is made.

\begin{theorem}\label{thm:three}
Suppose $h(a)=h(b)$, $a\ne b$, both inputs are cubefree, and they differ
at exactly three prime bases. After cancellation of unchanged full blocks,
their residual pair is $\{12,14\}$. Equivalently,
\[
 \{a,b\}=\{12d,14d\},\qquad (d,42)=1,
\]
where $d$ is cubefree.
\end{theorem}

\subsection{Two elementary divisibility lemmas}
Write $\Phi(x)=x^2+x+1$.
\begin{lemma}\label{lem:mixed}
If $p<r$ are primes, $r\mid\Phi(p)$ and $p\mid r+1$, then
$(p,r)=(2,7)$.
\end{lemma}
\begin{proof}
Write $r=\ell p-1$, $\ell\ge2$, and $c=\Phi(p)/r$.
Modulo $p$, $c\equiv-1$, so $c\ge p-1$. Hence
\[
 2p-1\le r\le p+2+\frac3{p-1}.
\]
This is impossible for $p\ge5$. For $p=3$ the only possible $r$ is $13$
and $3\nmid14$. For $p=2$ the only possible $r$ is $7$, which works.
\end{proof}

\begin{lemma}\label{lem:mutual}
If distinct primes $a,b$ satisfy $a\mid\Phi(b)$ and $b\mid\Phi(a)$,
then $a^2\nmid\Phi(b)$ and $b^2\nmid\Phi(a)$.
\end{lemma}
\begin{proof}
We first establish, for every coprime positive pair $x,y$ with the mutual
divisibilities, the identity
\begin{equation}\label{eq:five}
 x^2+y^2+x+y+1=5xy.
\end{equation}
For $1<x<y$ let $z=\Phi(x)/y$. Then $1\le z<x$: division by $y\ge x+1$
gives $z<x+1$, and $z=x$ would give $x\mid1$. Since $yz\equiv1\pmod x$,
the pair $(z,x)$ is coprime. Multiplying $y^2+y+1\equiv0\pmod x$ by
$z^2$ gives $z^2+z+1\equiv0\pmod x$; the other divisibility is the
definition of $z$. The value
\[
 \frac{x^2+y^2+x+y+1}{xy}=\frac{z+y+1}{x}
\]
is invariant under $(x,y)\mapsto(z,x)$. Descent reaches $(1,1)$ or $(1,3)$,
because the companion of $1$ divides $3$. Both give the value five, proving
\eqref{eq:five}.

If $a^2\mid\Phi(b)$, then \eqref{eq:five} gives
$\Phi(b)=a(5b-a-1)$ and $a\mid5b-1$. This rules out $a=5$; otherwise
substituting $5b\equiv1\pmod a$ into $25\Phi(b)\equiv0\pmod a$
gives $a\mid31$, hence $a=31$. But the quadratic for $b$ in
\eqref{eq:five} has discriminant
\[
 21\cdot31^2-14\cdot31-3=19744,
 \qquad 140^2<19744<141^2,
\]
so has no integer solution. Symmetry proves the other assertion.
\end{proof}

\subsection{Exhausting the possible local transitions}
Put
\[
 L_p=h(p)=p(p+1),\qquad S_p=h(p^2)=p^2\Phi(p),\qquad
 J_p=\frac{S_p}{L_p}=p^2+\frac p{p+1}.
\]
The local increments for states $0\leftrightarrow1$, $0\leftrightarrow2$,
and $1\leftrightarrow2$ are respectively $L_p,S_p,J_p$. All exceed one.
Moreover $p^2<J_p<p^2+1$, $J_p$ increases strictly with $p$, and its
reduced denominator is $p+1$, because $(p\Phi(p),p+1)=1$.

After canceling unchanged blocks, an equality on three distinct bases
equates one increment to the product of the other two. We sort by the number
of $J$ increments.

\textbf{Zero $J$ increments.} All indices are comparable. The Euler
product on three bases is at most $2(3/2)(5/4)=15/4<4$, so
Lemma~\ref{lem:comparable} excludes the equality.

\textbf{Two $J$ increments.} In the shape $J_p I_r=J_q$, where $I_r$ is
$L_r$ or $S_r$, denominators give $q+1\mid p+1$, hence $q\le p$.
But $I_r>1$ forces $J_q>J_p$ and $q>p$, a contradiction.
In the other shape $J_pJ_q=I_r$, both $p,q$ cannot be odd: the numerator
would be odd and the denominator even. If $p=2$, integrality of
\[
 J_2J_q=\frac{14q\Phi(q)}{3(q+1)}
\]
requires $q+1\mid14$, so $q=13$. The product is $793=13\cdot61$,
which is odd and squarefree; it is neither any $L_r$ nor any $S_r$.

\textbf{Three $J$ increments.} Suppose $J_pJ_q=J_r$.
Then $r+1\mid(p+1)(q+1)$. Also $r>pq$: for $r<pq$,
$J_r<r^2+1\le(pq-1)^2+1<p^2q^2<J_pJ_q$, while $pq$ is composite.
Since $(p-1)(q-1)\ge2$ for distinct primes,
\[
 (p+1)(q+1)\le2pq<2(r+1).
\]
The divisibility therefore forces $r+1=(p+1)(q+1)$, or $r=pq+p+q$.
But $J_r>r^2>(p^2+1)(q^2+1)>J_pJ_q$, a contradiction.

\textbf{One $J$ increment.} The noninteger $J_p$ cannot equal a product
of two integer increments. Thus $J_pI_q=I_r$, giving four cases.

\begin{enumerate}[label=\arabic*.]
\item $J_pL_q=L_r$. Comparison at $r=pq$ gives
$J_pq(q+1)>p^2q(q+1)>pq(pq+1)$, hence $r>pq$.
The cleared equation is
\[
 p\Phi(p)q(q+1)=(p+1)r(r+1).
\]
As $r>p,q+1$, it follows that $r\mid\Phi(p)$; also $p\mid r+1$.
Lemma~\ref{lem:mixed} gives $p=2,r=7$, then $q(q+1)=12$, so $q=3$.
This is $h(12)=h(14)=336$.

\item $J_pS_q=L_r$. Comparison at $r=pq^2$ gives
\[
 J_pq^2\Phi(q)>p^2q^2\Phi(q)>pq^2(pq^2+1),
\]
so $r>pq^2>\Phi(q)$ and $r>p$. The second strict inequality follows
from $2q^2>q^2+q+1$ for $q\ge2$.
The cleared equation again forces $r\mid\Phi(p)$ and $p\mid r+1$.
The mixed lemma yields $p=2,r=7$, contradicting $r>2q^2\ge18$ since
$q\ne2$.

\item $J_pS_q=S_r$. If all bases are odd, the two sides have respectively
negative and zero $2$-adic valuations. If $p=2$, the valuations are one
and zero. If $r=2$, the left side exceeds $28$. The remaining possibility
is $q=2$, giving $28J_p=S_r$ with $p,r$ odd. Its $2$-adic row requires
$v_2(p+1)=2$, and integrality requires $p+1\mid28$. Together these force
$p=3$, but $28J_3=273=3\cdot7\cdot13$ is squarefree and cannot be $S_r$.

\item $J_pL_q=S_r$. The case $p=2$ is excluded by parity. If $q=2$,
integrality requires $p+1\mid6$, so $p=5$; then $6J_5=155$ is squarefree.
If $r=2$, the left side exceeds $28$. Thus all three primes are odd, and
integrality gives $p+1\mid q(q+1)$.

If $p<q$, monotonicity in $S_x=J_xL_x$ gives $p<r<q$. As $q>p+1$,
we have the integer $c=(q+1)/(p+1)$ and
\[
 S_r=cpq\Phi(p),\qquad q\mid\Phi(r).
\]
If $r\mid c$, then $r\mid q+1$, contrary to the mixed lemma for the
odd pair $r<q$. Thus $r\nmid c$ and $r^2\mid\Phi(p)$, impossible
because $\Phi(p)<r^2$.

If $p>q$, monotonicity gives $q<r<p$. The divisibility
$p+1\mid q(q+1)$ and $p+1>q+1$ force $q\mid p+1$. Set
$c=(p+1)/q$, an even integer, and $t=(q+1)/c\in\NN$. Then
\[
 S_r=tp\Phi(p),\qquad t\le(q+1)/2<r.
\]
Hence $r^2\mid\Phi(p)$ and $p\mid\Phi(r)$, contradicting
Lemma~\ref{lem:mutual}.
\end{enumerate}
The cases exhaust all increments. Restoring the unchanged full blocks proves
Theorem~\ref{thm:three}. \qed

Combined with the separately audited fixed-two-prime theorem from the handoff,
this implies that fixing the input exponent at $2$ makes distinct cubefree
preimages differ at at least four bases. Minimum distance four still permits
exponentially many abstract exponent patterns. No global injectivity on
cubefree inputs follows.

\section{Finite verification and the remaining obstruction}
\subsection{Fresh exact checks}
The new standard-library verifier uses complete divisor enumeration for the
following unrestricted fibers; $k\mid h(k)$ makes each candidate list exhaustive.
\begin{center}\small
\begin{tabular}{@{}r r r@{}}
\toprule Target & Divisors tested & Fiber size\\\midrule
$106345074572730040320$ & $7424$ & $7$\\
$7089671638182002688000$ & $12288$ & $6$\\
$1826980530660612389572800675840$ & $23552$ & $9$\\
\bottomrule
\end{tabular}
\end{center}
An independent local-state enumeration also reproduces the cap-two fibers
of the supplied large target $M$ and $336M$, of sizes two and four, and the
singleton cap-two fiber of $504654433920$. Its pruning only removes states
that violate necessary valuation-row bounds; every final row and returned
input is checked exactly. Resource exhaustion raises an error rather than
reporting completeness. These runs completed in 44, 51 and 9 search nodes.

A separate verifier checks the large odd cubefree equality, both denominator
normalizations, its 16-label obstruction, and all four minimal exchanges of
the seven-member fiber, including proper subset tests. Another finite
regression examines all three-base transition types on 669 primes below
5000: 2,011,014 distinct-base increment-pair comparisons return only the
residual pair $\{12,14\}$. This is a regression check of the proof, not its
justification for primes outside the tested range.

The full historical census through $10^{16}$ and every large stored
certificate were not rerun in this audit. Their historical output is not
being promoted to new independent verification. Exact inputs, output lists,
and code are in the accompanying audit bundle.

\subsection{Why the new packing statements still do not finish the proof}
Let $\Gamma(k_0)$ be the minimal exchanges, $g=|\Gamma(k_0)|$, $\nu$ the
largest disjoint collection, and $Z$ the number of disjoint collections.
The proved decomposition yields only
\[
 \max\{g+1,2^\nu\}\le f_E(N)\le Z
 \le\sum_{j=0}^{\nu}\binom gj.
\]
An upper bound on $\nu$ is not the estimate $f_E(N)\le2^\nu$. In the
actual seven-member example, $\nu=2$ and $f(N)=7>4$. Many exchanges can
overlap and still give many alternatives. The denominator budget controls
bounded-denominator disjoint collections; it leaves both large denominators
and the number of overlapping alternatives uncontrolled.

Likewise the valid inequality, for $\theta>0$,
\[
 f_E(N)\le N^\theta
 \prod_{\gamma\in\Gamma(k_0)}(1+H_\gamma^{-\theta})
\]
does not estimate its right-hand side at the required scale. Enumerating the
entire fiber first to construct $\Gamma(k_0)$ cannot serve as a uniform
upper-bound proof.

\subsection{Further routes examined and what they would require}
\begin{enumerate}
\item \textbf{Denominator stratification.} The new budget makes the bounded
denominator part of a disjoint collection small. A successful argument would
also need to compress the large-denominator exchanges and control overlaps,
using arithmetic properties not supplied by the decomposition.
\item \textbf{Small-support classification.} Three changed bases at cap two
are now settled. The next step would need control of larger minimal exchanges
uniformly in support size and prime magnitude. A finite classification at
each of a few support sizes does not give that control.
\item \textbf{Pairwise reconstruction.} Proposition~\ref{prop:pairwise}
is an exact sufficient criterion. A useful pairwise denominator bound for all
fixed-cap fibers is unproved; estimates only against one reference do not
establish it.
\item \textbf{Cyclotomic and entropy methods.} Primitive divisors can be
shared by many bases; a fixed-length average product theorem need not control
one worst-case growing product. A uniformly small arithmetic encoding is
still needed before chain, graph, or integer-kernel machinery can close the gap.
\end{enumerate}

The exact outstanding statement remains \eqref{eq:fixedcap} for every fixed
$E$. The current work proves neither that statement nor a counterexample to it.
Accordingly, the general coefficient remains $C_1>0$ in the reviewed record.
The complete proofs of the additional restricted theorems above should not be
presented as a complete proof of Erd\H{o}s 1060.

\begin{thebibliography}{99}
\small
\bibitem{bloom} Thomas Bloom, Erd\H{o}s Problems, Problem 1060.
Indexed listing inspected 7 September 2026; direct retrieval blocked.
\url{https://www.erdosproblems.com/1060}.
\bibitem{kominers} Scott Duke Kominers, \emph{On the Number of Solutions of
$k\sigma(k)=n$}, working paper, 2026.
\url{https://www.scottkom.com/assets/articles/Kominers_ksigmak.pdf}.
\bibitem{np} Passawan Noppakaew and Prapanpong Pongsriiam,
\emph{Product of Some Polynomials and Arithmetic Functions},
Journal of Integer Sequences 26 (2023), Article 23.9.1.
\url{https://cs.uwaterloo.ca/journals/JIS/VOL26/Pongsriiam/pong43.pdf}.
\bibitem{guy} Richard K. Guy, \emph{Unsolved Problems in Number Theory},
third edition, Springer, 2004, B11, pp. 96--97. The corresponding source
pages were inspected in the available sample.
\url{https://doi.org/10.1007/978-0-387-26677-0}.
\bibitem{erdos} Paul Erd\H{o}s, \emph{Remarks on number theory II
Some problems on the $\sigma$ function}, 1959, especially Theorem 2.
\url{https://www.renyi.hu/~p_erdos/1959-21.pdf}.
\bibitem{pollack} Paul Pollack, \emph{Remarks on fibers of the sum-of-divisors
function}, especially Theorem 2.
\url{https://www.pollack-math.net/preimages-maier3.pdf}.
\bibitem{yamada} Tomohiro Yamada, \emph{Multiplicative structures of values of
the sum-of-divisors function}, arXiv:math/0512175 (2005), Theorem 1.4.
\url{https://arxiv.org/abs/math/0512175}.
\bibitem{wx} Victor Y. Wang and Max Wenqiang Xu,
\emph{Paucity phenomena for polynomial products},
Bulletin of the London Mathematical Society 56 (2024), 2718--2726.
\url{https://arxiv.org/abs/2211.02908}.
\bibitem{ess} J.-H. Evertse, H. P. Schlickewei and W. M. Schmidt,
\emph{Linear equations in variables which lie in a multiplicative group},
Annals of Mathematics 155 (2002), 807--836.
\url{https://annals.math.princeton.edu/2002/155-3/p04}.
\bibitem{mills} W. H. Mills, \emph{A system of quadratic Diophantine equations},
Pacific Journal of Mathematics 3 (1953), 209--220.
\url{https://msp.org/pjm/1953/3-1/pjm-v3-n1-p15-p.pdf}.
\bibitem{bibby} Sean Bibby, Pieter Vyncke and Joshua Zelinsky,
\emph{On the Third Largest Prime Divisor of an Odd Perfect Number},
INTEGERS 21 (2021), A115.
\url{https://arxiv.org/abs/1908.09420}.
\bibitem{richard} Adrien Richard, \emph{Positive circuits and maximal number
of fixed points in discrete dynamical systems}, Discrete Applied Mathematics
157 (2009), 3281--3288.
\url{https://arxiv.org/abs/0807.4229}.
\bibitem{graver} Dinh Van Le and Tim R\"omer, \emph{Equivariant lattice bases},
arXiv:2309.07246, Definition 2.2 and Remark 2.3.
\url{https://arxiv.org/abs/2309.07246}.
\end{thebibliography}
\end{document}
